O mercado de adquirência no Brasil, em constante expansão e com elevado volume de transações, exige da logística uma resposta estratégica para gerenciar o ciclo de vida das maquinetas POS. Esse cenário, que incorpora a logística direta e a LR, caracteriza-se como uma CSCF, apresentando desafios críticos de localização de facilidades e otimização de custos operacionais. Diante disso, o cerne desta dissertação foi a proposta e aplicação de um modelo PLIM (Capítulo \ref{chapter:def_prob_mod}). O objetivo fundamental era a minimização do custo logístico operacional total da cadeia, abrangendo as decisões de alocação de CDs e PAs, os fluxos de movimentação e os níveis de estoque ao longo de um horizonte de um ano.

Os experimentos computacionais, executados com o solver Gurobi (Capítulo \ref{chapter:exp_comp}), foram determinantes para validar a capacidade do modelo PLIM em alcançar a melhor solução factível global encontrada ($gap \leq 5\%$) para a CSCF em um contexto de alta complexidade. A investigação dos resultados demonstrou a eficácia em gerar a configuração de custo mínimo para a cadeia, atendendo plenamente à função objetivo. Além disso, o modelo de otimização também serviu como uma ferramenta de diagnóstico, pois a análise aprofundada da utilização das restrições e dos fluxos internos expôs os pontos de estrangulamento da operação. É com esta dupla perspectiva: otimização de custos e exposição de ineficiências operacionais, que a análise detalhada das conclusões se inicia.

O modelo PLIM demonstrou alta eficácia, atingindo um resultado ótimo global e expondo a ineficiência operacional da cadeia logística. O principal gargalo identificado foi o fabricante, cujos parâmetros, assim como os dos PAs, estavam superestimados. Em todas as instâncias, o fabricante operou em 100\% de sua capacidade, indicando que o fornecimento e o recebimento de materiais pelo CD estavam abaixo do necessário.

Por sua vez, o CD executou seu papel com eficiência, suprindo a fragilidade do fabricante. No período analisado, sua atuação foi suficiente; entretanto, em um horizonte temporal maior, essa condição pode gerar problemas. Isso porque o CD foi obrigado a acumular itens para envio aos PAs e, devido à Restrição \ref{re9_3}, que estabelece um estoque mínimo por falta de uma política de estoque definida, a capacidade de armazenamento foi utilizada em 100\%.

Com relação aos PAs, observa-se subutilização em todas as restrições, indicando que o modelo optou por priorizar a otimização de custos relacionados à abertura de PAs em detrimento das movimentações. 

Financeiramente, o modelo demonstrou excelente desempenho. No entanto, trata-se de um caso que merece análises mais aprofundadas, uma vez que, em algumas situações, pode resultar em rotas inviáveis e operacionalmente ineficientes para um sistema real.

Para análises futuras, destacam-se pontos relevantes que abrangem desde aspectos de pensamento computacional até questões logísticas de gestão de estoques.

Na dimensão computacional, recomenda-se desenvolver um modelo que incorpore uma solução heurística. Esta adaptação permitiria a aplicação em escala nacional, estendendo-se além do escopo estadual de São Paulo adotado no estudo atual. Métodos como \sigla{RF}{\emph{Relax-and-Fix}} ou \sigla{FO}{\emph{Fix-and-Optimize}} e outras abordagens de decomposição temporal podem ser exploradas. O uso de técnicas de decomposição é fundamental para resolução integral de problemas de localização com múltiplos períodos e custos dinâmicos. Seguindo esta linha de estudo, \citeonline{HuacchaNeyra2017} utiliza a heurística híbrida RF-FO para a otimização de dimensionamento de lotes, que mostrou-se mais rápida com soluções de melhor qualidade do que a resolução via solver CPLEX.

No âmbito logístico, embora o modelo em questão tenha alcançado a otimização de custos, torna-se imperativa a adoção de políticas de estoque que possibilitem a construção de cenários mais aderentes à realidade operacional. Isso inclui a previsão de estoques de segurança em todos os níveis da rede e a ampliação de variáveis e parâmetros associados aos fabricantes. Nesse contexto, \citeonline{Sassaki2022Politica} propõe um estudo sobre políticas de estoque multinível orientadas à maximização da eficiência na CSCF de uma empresa de adquirência. Tal expansão de escopo visa mitigar a formação de gargalos, evitar a superestimação de capacidades e fomentar uma gestão de recursos mais equilibrada e eficaz.

Por fim, o modelo mostrou-se matematicamente eficiente e, em um contexto real e prático, pode servir como um poderoso balizador para a estruturação de uma cadeia logística como a estudada. Ressalta-se, contudo, a necessidade de reflexão sobre situações que podem gerar ineficiências logísticas em função da distância entre os PAs e as ZDs.
